\documentclass[10pt,a4paper]{amsart}
\usepackage[margin=19mm]{geometry}
\usepackage{amsmath,amssymb,mathtools}
\usepackage[scr=rsfs,cal=euler]{mathalfa}
\usepackage{xfrac}
\usepackage[unicode,hidelinks,bookmarks=false]{hyperref}
\newcommand{\ie}{i.e.\ }
\newcommand{\cf}{cf.\ }
\newcommand{\resp}{respectively\ }
\newcommand{\Subsection}{\subsection}
\providecommand{\nobreakdash}{\nobreak\hbox{-}}
\setlength{\emergencystretch}{2em}
\multlinegap0pt
\usepackage{amssymb}
\usepackage[normalem]{ulem}
\usepackage{xcolor}
\newtheorem{theorem}{Theorem}
\title{Extending recent work of Nath, Saikia, and Sarma on $k$-tuple $\ell$-regular Partitions}
\author{Bishnu Paudel, James A. Sellers,, Haiyang Wang}
\date{Source: arXiv:2503.12583v1 (CC BY 4.0)}
\begin{document}
\maketitle

\noindent\emph{Source and licence.} Extending recent work of Nath, Saikia, and Sarma on $k$-tuple $\ell$-regular Partitions. Version arXiv:2503.12583v1 (CC BY 4.0), distributed by arXiv under the Creative Commons Attribution 4.0 International licence, http://creativecommons.org/licenses/by/4.0/, as recorded in the arXiv licence metadata for this exact version and shown on the abstract page https://arxiv.org/abs/2503.12583v1. Full licence notice: \texttt{licenses/arxiv-2503.12583-CC-BY-4.0.txt}.\par\smallskip

\noindent\textit{Editorial scope.} Counted unit: the article's theorem thm.p357mod8 (source0.tex lines 189--194). The article's complete original proof of it is delivered with it (source0.tex lines 386--422, one \texttt{proof} environment of 37 lines, which the article places away from the statement). No mathematics is written, repaired or re-derived: the statement and the proof are the article's own lines, in the article's own notation, and no retained line is substituted or re-flowed.  Retained uncounted as the source-local context the proof needs: lines 234--236; lines 242--244; lines 251--253; lines 259--261; lines 283--285; lines 294--296; lines 303--304.  Nothing was omitted from any retained span, so the package carries no omission record.  No retained line contains a citation, so the wrapper carries no bibliography and no bibliography entry.  No retained line contains a figure environment, an image-inclusion command or drawing code, so no image is delivered and the package declares no asset dependency.  Editorial formatting only, in addition: an AMS \texttt{amsart} preamble that restates the article's own declarations and macros used by the retained lines, namely the environments \texttt{theorem} on separate counters, together with the standard packages those lines need.  The retained environments print in this document as Theorem 1 in order of appearance, whereas the article prints the same items with the numbering of its own full document.  The complete native source of the article as distributed in the arXiv e-print for this exact version is bundled unchanged as \texttt{sources/174862/source0.tex}.
    \begin{equation}\label{f1}
        f_1=\sum_{m\in\mathbb{Z}}(-1)^mq^{m(3m-1)/2}.
    \end{equation}

\begin{equation}\label{f13}
    f_1^3=\sum_{m\ge 0}(-1)^m(2m+1)q^{m(m+1)/2}.
\end{equation}

  \begin{equation}\label{f15/f22}
        \frac{f_1^5}{f_2^2}=\sum_{m\in\mathbb{Z}}(6m+1)q^{m(3m+1)/2}.
    \end{equation}  

\begin{equation}\label{f1-q}
    (-q;-q)_\infty=\frac{f_2^3}{f_1 f_4}.
\end{equation}

\begin{equation}\label{f22/f4}
    \frac{f_2^2}{f_4}=1+2\sum\limits_{n\ge 1}(-1)^nq^{2n^2}.
\end{equation}

    \begin{equation}\label{modp^k}
        f_l^{p^k}\equiv f_{lp}^{p^{k-1}}\pmod{p^k}.
    \end{equation}

\begin{equation}\label{modp}
T_{\ell,p}(pn+r)\equiv0\pmod{p}.\end{equation}

\begin{theorem}\label{thm.p357mod8}
	Let $p\equiv3, 5 \text{ or } 7\mod 8$ be a prime. Then, for $n,\alpha\geq 0$ with $p\nmid n$, we have $$T_2\left(p^{2\alpha+1}n+\frac{p^{2\alpha+2}-1}{8}\right)\equiv0\pmod{8}.$$
In addition, if $3\nmid n$, then 
\begin{equation}\label{mod24}
T_2\left(p^{2\alpha+1}n+\frac{p^{2\alpha+2}-1}{8}\right)\equiv0\pmod{24}.\end{equation}
\end{theorem}

\begin{proof}[Proof (of Theorem \ref{thm.p357mod8})]
 Thanks to \eqref{modp^k}, we have 
\begin{equation}\label{mod8}
    \sum_{n\geq0}T_2(n)q^n=\frac{f_2^3}{f_1^3}=\frac{f_2^4}{f_1^3f_2}\equiv\frac{f_1^5}{f_2^2}f_2 \pmod8.
\end{equation}
Using \eqref{f1} and \eqref{f15/f22} in \eqref{mod8}, we obtain 
\begin{align*}
    \sum_{n\geq0}T_2(n)q^n\equiv \sum_{m\in\mathbb{Z}}\sum_{k\in\mathbb{Z}}(-)^k(6m+1)q^{m(3m+1)/2+k(3k-1)} \pmod 8.
\end{align*}
We now need to check whether $l:=p^{2\alpha+1}n+\cfrac{p^{2\alpha+2}-1}{8}$ can be represented as $\cfrac{m(3m+1)}{2}+k(3k-1)$. Equivalently, we check whether $$24l+3=(6m+1)^2+2(6k-1)^2.$$ 
Let $\nu_p(N)$ be the highest power of $p$ dividing $N$. We first consider primes $p\equiv5\text{ or }7\mod 8$. Then, we have $\left(\cfrac{-2}{p}\right)=-1$. So, if $N=x^2+2y^2$, then $2|\nu_p(N)$. Also, for $p\nmid n$, we have $\nu_p(24l+3)=2\alpha+1$. Thus, $24l+3$ cannot be of the form $x^2+2y^2$. Hence, for $p\equiv5 \text{ or } 7 \mod 8$,
$$T_2\left(p^{2\alpha+1}n+\frac{p^{2\alpha+2}-1}{8}\right)\equiv 0\pmod{8}.$$

The above proof only deals with the cases $p\equiv5 \text{ or } 7 \mod 8$, which still leaves us with the case $p\equiv3 \mod 8$.  This requires a different idea.  In this case, we replace 
$q$ by $-q$ and then, using \eqref{f1-q} and \eqref{modp^k}, we obtain 
\begin{equation} \label{mod82}
    \sum_{n\geq0}T_2(n)(-q)^n =f_2^3\left(\frac{f_1f_4}{f_2^3}\right)^3=\frac{f_1^3f_4^3f_2^2}{f_2^8}\equiv \frac{f_1^3f_4^3f_2^2}{f_4^4}= f_1^3\frac{f_2^2}{f_4} \pmod{8}.
\end{equation}
Substituting \eqref{f13} and \eqref{f22/f4} in \eqref{mod82} gives
\begin{align*}
   \sum_{n\geq0}T_2(n)(-q)^n\equiv \sum_{m\in\mathbb{Z}}(-1)^m(2m+1)q^{m(m+1)/2}\left(1+2\sum_{k\geq 1}(-1)^kq^{2k^2}\right) \pmod{8}.
\end{align*}
Now we need to check whether we have 
$$l=\frac{m(m+1)}{2}+2k^2,$$
that is, $$8l+1=(2m+1)^2+(4k)^2.$$
For primes $p\equiv 3\mod 4$, we note that $\left(\cfrac{-1}{p}\right)=-1$. So, for any $N=x^2+y^2$, we have $2|\nu_p(N)$. However, for $p\nmid n$,  $\nu_p(8l+1)=2\alpha+1$. This completes the proof of the congruences modulo 8.

For the second part, it suffices to show that $T_2(l)\equiv 0\pmod{3}$. For this, we check whether $3|l$ (in order to apply \eqref{modp}). Since $8^{-1}\equiv-1 \pmod 3$, we have 
\[l=p^{2\alpha+1}n+\cfrac{p^{2\alpha+2}-1}{8}\equiv
\begin{cases}n \pmod{3} &\text{ if } p\equiv1\pmod{3},\\
-n \pmod{3} & \text{ if } p\equiv -1\pmod {3},\\
1 \pmod{3} & \text{ if } p=3.
\end{cases}
\]
Since $3\nmid n$, we have $3\nmid l$ and, by \eqref{modp}, $T_2(l)\equiv0\pmod 3$.

\end{proof}

\end{document}
